\documentclass[reqno, answers]{exam}

\usepackage{amssymb, amsmath, amsthm, stmaryrd}
\usepackage{fullpage, parskip}
\usepackage{enumitem}
\usepackage{tikz}
\usepackage{ulem}
\usepackage{hyperref}

\title{Introduction to Computational Math Spring 2026 \\ Homework 02}
\date{Due: Friday, Feb 06, 2025, 11:59 PM}
\author{}

\begin{document}
\maketitle
\thispagestyle{empty}

Textbook = Driscoll and Braun (\url{https://fncbook.com/})

\begin{enumerate}

\item Textbook Exercise 1.2.3

\begin{solution}

\textbf{You can do all of these using L'Hôpital's rule, but I've given solutions here using Taylor series expansions. Both methods are correct.}

The relative condition number of a function $f(x)$ is defined as:
$$\kappa_f(x) = \left|\frac{xf'(x)}{f(x)}\right|$$

We find $\kappa_f(x) \to \infty$ when either $|xf'(x)| \to \infty$ or $f(x) \to 0$ (or both).

\textbf{(a)} $f(x) = \tanh(x)$

We have $f'(x) = \operatorname{sech}^2(x) = 1 - \tanh^2(x)$.

Using the exponential form:
$$f(x) = \tanh(x) = \frac{e^{2x} - 1}{e^{2x} + 1}, \quad f'(x) = \frac{4e^{2x}}{(e^{2x} + 1)^2}$$

Therefore:
$$\kappa_f(x) = \left|\frac{x \cdot \frac{4e^{2x}}{(e^{2x} + 1)^2}}{\frac{e^{2x} - 1}{e^{2x} + 1}}\right| = \left|\frac{4xe^{2x}}{(e^{2x} + 1)(e^{2x} - 1)}\right|$$

\textit{Finding where $\kappa_f(x) \to \infty$:}

- At $x = 0$: Both $e^{2x} - 1 \to 0$ and $x \to 0$. Using L'Hôpital's rule or Taylor series with $e^{2x} - 1 \approx 2x$ and $e^{2x} + 1 \approx 2$ as $x \to 0$:
$$\lim_{x \to 0} \kappa_f(x) = \lim_{x \to 0} \left|\frac{4x \cdot 1}{2 \cdot 2x}\right| = 1$$
So $\kappa_f(0) = 1$ is finite.

- As $x \to +\infty$: $\tanh(x) \to 1$, so $f(x) \to 1$ (not zero), and $\kappa_f(x) \to 0$.

- As $x \to -\infty$: $\tanh(x) \to -1$, so $f(x) \to -1$ (not zero), and $\kappa_f(x) \to 0$.

Therefore, $\boxed{\kappa_f(x) \text{ never approaches } \infty}$.

\textbf{(b)} $f(x) = \frac{e^x - 1}{x}$

Using the quotient rule: $f'(x) = \frac{xe^x - (e^x - 1)}{x^2}$

$$\kappa_f(x) = \left|\frac{x \cdot \frac{xe^x - e^x + 1}{x^2}}{\frac{e^x - 1}{x}}\right| = \left|\frac{xe^x - e^x + 1}{(e^x - 1)}\right|$$

\textit{Finding where $\kappa_f(x) \to \infty$:}

- At $x = 0$: Using Taylor series:
$$e^x - 1 \approx x + \frac{x^2}{2}, \quad xe^x - e^x + 1 \approx \frac{x^2}{2}$$
$$\kappa_f(x) \approx \left|\frac{x^2/2}{x^2}\right| = \frac{1}{2}$$
So $\kappa_f(0) = \frac{1}{2}$ is finite.


- As $x \to +\infty$: 
$$\kappa_f(x) = \left|\frac{xe^x - e^x + 1}{e^x - 1}\right| = \left|\frac{e^x(x - 1) + 1}{e^x - 1}\right| \to |x - 1| \to \infty$$

- As $x \to -\infty$: $e^x \to 0$, so 
$$\kappa_f(x) = \left|\frac{0 - 0 + 1}{0 - 1}\right| = 1$$

Therefore, $\boxed{\kappa_f(x) \to \infty \text{ as } x \to +\infty}$.

\textbf{(c)} $f(x) = \frac{1 - \cos(x)}{x}$

Using the quotient rule: $f'(x) = \frac{x\sin(x) - (1-\cos(x))}{x^2}$

$$\kappa_f(x) = \left|\frac{xf'(x)}{f(x)}\right| = \left|\frac{x \cdot \frac{x\sin(x) - (1-\cos(x))}{x^2}}{\frac{1-\cos(x)}{x}}\right| = \left|\frac{x\sin(x) - (1-\cos(x))}{1-\cos(x)}\right|$$

\textit{Finding where $\kappa_f(x) \to \infty$:}

- At $x = 0$: Using Taylor series:
$$\sin(x) = x - \frac{x^3}{6} + \cdots, \quad 1 - \cos(x) = \frac{x^2}{2} - \frac{x^4}{24} + \cdots$$
$$x\sin(x) = x^2 - \frac{x^4}{6} + \cdots$$
$$x\sin(x) - (1-\cos(x)) = x^2 - \frac{x^4}{6} - \frac{x^2}{2} + \frac{x^4}{24} + \cdots = \frac{x^2}{2} - \frac{x^4}{8} + \cdots$$

$$\kappa_f(x) = \left|\frac{\frac{x^2}{2} - \frac{x^4}{8} + \cdots}{\frac{x^2}{2} - \frac{x^4}{24} + \cdots}\right| = \left|\frac{x^2(\frac{1}{2} - \frac{x^2}{8} + \cdots)}{x^2(\frac{1}{2} - \frac{x^4}{24} + \cdots)}\right| = \left|\frac{\frac{1}{2} - \frac{x^2}{8} + \cdots}{\frac{1}{2} - \frac{x^2}{24} + \cdots}\right| \to 1$$

So $\kappa_f(0) = 1$ is finite.

- At $x = 2\pi n$ for nonzero integer $n$: We have $1 - \cos(2\pi n) = 0$, so the denominator goes to zero. Near $x = 2\pi n + h$:
$$1 - \cos(2\pi n + h) = 1 - \cos(h) \approx \frac{h^2}{2}$$
$$\sin(2\pi n + h) = \sin(h) \approx h$$
$$(2\pi n + h)\sin(2\pi n + h) \approx 2\pi n \cdot h + h^2 = h(2\pi n + h)$$
$$x\sin(x) - (1-\cos(x)) \approx h(2\pi n + h) - \frac{h^2}{2} = 2\pi nh + \frac{h^2}{2}$$

$$\kappa_f(x) \approx \left|\frac{2\pi nh + \frac{h^2}{2}}{\frac{h^2}{2}}\right| = \left|\frac{4\pi n}{h} + 1\right| \to \infty$$

Therefore, $\boxed{\kappa_f(x) \to \infty \text{ as } x \to 2\pi n \text{ for all nonzero integers } n}$.
\end{solution}

\item Textbook Exercise 1.2.6
\begin{solution}

Consider the quadratic equation $ax^2 + bx + c = 0$ with roots $r_1$ and $r_2$.

We want to find the condition number of a root $r$ with respect to $b$:
$$\kappa = \left|\frac{b}{r} \cdot \frac{\partial r}{\partial b}\right|$$


Starting with the equation $ax^2 + bx + c = 0$, where $x = r$ is a root that depends on $b$.

Differentiating implicitly with respect to $b$:
$$a \cdot 2r \cdot \frac{\partial r}{\partial b} + b \cdot \frac{\partial r}{\partial b} + r \cdot 1 + 0 = 0$$

$$2ar \cdot \frac{\partial r}{\partial b} + b \cdot \frac{\partial r}{\partial b} + r = 0$$

$$(2ar + b) \frac{\partial r}{\partial b} = -r$$

$$\frac{\partial r}{\partial b} = \frac{-r}{2ar + b}$$


$$\kappa = \left|\frac{b}{r} \cdot \frac{\partial r}{\partial b}\right| = \left|\frac{b}{r} \cdot \frac{-r}{2ar + b}\right| = \left|\frac{-b}{2ar + b}\right| = \left|\frac{b}{2ar + b}\right|$$


Plugging in $r_1 = \frac{-b + \sqrt{b^2 - 4ac}}{2a}$ this simplifies to:

$$\kappa = \left|\frac{b}{\sqrt{b^2 - 4ac}}\right|$$

Finally, plugging in $\sqrt{b^2 - 4ac} = a(r_1 - r_2)$ gives the result:

$$\kappa = \left|\frac{b}{a(r_1 - r_2)}\right|$$

The other root has the same condition number with respect to $b$.

\end{solution}

\item Textbook Exercise 1.2.8

\begin{solution}
The relative condition number is:
$$\kappa_r(a_k) = \left|\frac{a_k}{r} \cdot \frac{\partial r}{\partial a_k}\right|$$


Since $r$ is a root, we have:
$$p(r) = a_n r^n + a_{n-1} r^{n-1} + \cdots + a_k r^k + \cdots + a_1 r + a_0 = 0$$

Treating $r$ as a function of $a_k$ and differentiating implicitly with respect to $a_k$:
$$a_n \cdot n r^{n-1} \frac{\partial r}{\partial a_k} + a_{n-1} \cdot (n-1) r^{n-2} \frac{\partial r}{\partial a_k} + \cdots + a_k \cdot k r^{k-1} \frac{\partial r}{\partial a_k} + r^k + \cdots + a_1 \frac{\partial r}{\partial a_k} = 0$$

Notice that all terms containing $\frac{\partial r}{\partial a_k}$ sum to:
$$\left(a_n \cdot n r^{n-1} + a_{n-1} \cdot (n-1) r^{n-2} + \cdots + a_1\right) \frac{\partial r}{\partial a_k} = p'(r) \frac{\partial r}{\partial a_k}$$

The only term without $\frac{\partial r}{\partial a_k}$ is $r^k$ (from differentiating $a_k r^k$ with respect to $a_k$).

Therefore:
$$p'(r) \frac{\partial r}{\partial a_k} + r^k = 0$$

$$\frac{\partial r}{\partial a_k} = -\frac{r^k}{p'(r)}$$


$$\kappa_r(a_k) = \left|\frac{a_k}{r} \cdot \frac{\partial r}{\partial a_k}\right| = \left|\frac{a_k}{r} \cdot \left(-\frac{r^k}{p'(r)}\right)\right|$$

$$= \left|\frac{-a_k r^k}{r \cdot p'(r)}\right| = \left|\frac{-a_k r^{k-1}}{p'(r)}\right|$$

$$= \boxed{\left|\frac{a_k r^{k-1}}{p'(r)}\right|}$$

\end{solution}


\item \sout{Textbook Exercise 1.4.1}

\end{enumerate}

\end{document}